%% ============================================================
%%  CV – Damien MAÎTRE  |  Versione italiana
%%  Engine : XeLaTeX  |  Font : Calibri (or fallback Linux)
%% ============================================================
\documentclass[a4paper,9pt]{extarticle}

% ── Packages ─────────────────────────────────────────────────
\usepackage[T1]{fontenc}
\usepackage[utf8]{inputenc}
\usepackage{lmodern}
\usepackage[left=1.8cm, right=1.8cm, top=1.2cm, bottom=1.2cm]{geometry}
\usepackage{xcolor}
\usepackage{fontawesome5}
\usepackage{titlesec}
\usepackage{enumitem}
\usepackage{hyperref}
\usepackage{array}
\usepackage{tabularx}
\usepackage{booktabs}
\usepackage{microtype}
\setlength{\parskip}{1pt}
\usepackage{multicol}

% ── Colours ──────────────────────────────────────────────────
\definecolor{accent}{HTML}{1A3C6E}
\definecolor{light}{HTML}{4A7CB5}
\definecolor{muted}{HTML}{555555}
\definecolor{rule}{HTML}{CCCCCC}

% ── Hyperlinks ───────────────────────────────────────────────
\hypersetup{
  colorlinks=true,
  urlcolor=light,
  linkcolor=accent
}

% ── No page number ───────────────────────────────────────────
\pagestyle{empty}

% ── Section titles ───────────────────────────────────────────
\titleformat{\section}
  {\large\bfseries\color{accent}}
  {}{0em}{}
  [\vspace{1pt}\color{rule}\hrule\vspace{2pt}]

\titlespacing{\section}{0pt}{5pt}{2pt}

% ── List settings ────────────────────────────────────────────
\setlist[itemize]{
  leftmargin=1.4em,
  itemsep=0pt,
  topsep=1pt,
  parsep=0pt
}
\renewcommand{\labelitemi}{\textcolor{light}{\small$\bullet$}}

% ── Entry command: \entry{dates}{title}{subtitle}{bullets} ───
\newcommand{\entry}[4]{%
  \noindent
  \begin{tabularx}{\linewidth}{@{}p{2.6cm}@{\hspace{6pt}}X@{}}
    \small\textcolor{muted}{#1} &
    \textbf{#2} \newline
    \textit{\small\textcolor{muted}{#3}}
    #4
  \end{tabularx}
  \vspace{1pt}
}

\newcommand{\wideentry}[3]{%
  \noindent\textbf{#1}\\
  \textit{\small\textcolor{muted}{#2}}\\
  \small #3\par
  \vspace{1pt}
}

% ── Skill row: \skillrow{category}{items} ────────────────────
\newcommand{\skillrow}[2]{%
  \noindent
  \begin{tabularx}{\linewidth}{@{}p{3.0cm}@{\hspace{4pt}}X@{}}
    \small\textbf{\textcolor{accent}{#1}} & \small #2
  \end{tabularx}\par\vspace{1pt}
}

%% ============================================================
\begin{document}

% ── NAME & CONTACT ───────────────────────────────────────────
\begin{center}
  {\LARGE\bfseries\color{accent} Damien MAÎTRE}\\[2pt]
  {\normalsize\textit{\textcolor{muted}{Laurea Magistrale in Mobility Engineering — Infrastrutture, Meccanica \& Metodi Numerici}}}\\[4pt]
  \small
  \textcolor{muted}{
    \faEnvelope\ \href{mailto:damien.maitre@enpc.fr}{damien.maitre@enpc.fr}
    \quad\textbf{·}\quad
    \faPhone\ +33 7 83 41 50 90
    \quad\textbf{·}\quad
    \faLinkedin\ \href{https://linkedin.com/in/damienmaitre}{linkedin.com/in/damienmaitre}
    \quad\textbf{·}\quad
    \faGithub\ \href{https://github.com/dammaitre}{github.com/dammaitre}
  }
\end{center}

\vspace{2pt}
{\color{accent}\hrule height 1.2pt}
\vspace{4pt}

% ── ABOUT ME ─────────────────────────────────────────────────
\section{PROFILO}

\noindent\small
Sono uno studente all'ultimo anno di ingegneria presso l'École nationale des ponts et chaussées e il Politecnico di Milano, con esperienza concreta nil cantiere infrastrutturale. Sono alla ricerca di un tirocinio di fine studi nell'ambito della gestione o progettazione di progetti infrastrutturali, dell'ingegneria delle infrastrutture, oppure della progettazione, gestione operativa o ricerca nei servizi ferroviari e di mobilità. Attualmente vivo a Milano e sono disponibile a iniziare a febbraio o marzo 2027. \par

% ── EDUCATION ────────────────────────────────────────────────
\section{FORMAZIONE}

\entry{settembre 2026 \newline febbraio 2027}
  {Laurea Magistrale in Mobility Engineering}
  {Programma di scambio presso il Politecnico di Milano}
  {\newline Economia internazionale, esercizio ferroviario, calcolo ad alte prestazioni, finanza delle infrastrutture \& gestione delle infrastrutture.}

\entry{settembre 2023 \newline giugno 2027}
  {Laurea Magistrale in Ingegneria — Fluidodinamica, Meccanica delle strutture, Prototipazione}
  {École nationale des ponts et chaussées (ENPC), Parigi}
  {\newline Meccanica delle strutture e progettazione di grandi strutture, metodo agli elementi finiti (FEM) applicato all'analisi strutturale, materiali compositi, metallurgia, modulo completo di fluidodinamica, HPC, combustione.
  \begin{itemize}
    \item \textbf{Strumenti:} OpenFOAM, StarCCM+, Abaqus, OpenSeesPy, gmsh
    \item \textbf{Progetti:} simulazione e prototipo di un aliante subacqueo (sea-glider); progettazione e ottimizzazione di un velivolo ad ala chiusa; progettazione, verifica e costruzione di un ponte sospeso in legno.
  \end{itemize}}

\entry{settembre 2021 \newline giugno 2023}
  {Laurea triennale — \textit{Classes préparatoires}}
  {Lycée Kléber, Strasburgo}
  {\newline Programma intensivo di matematica, fisica e scienze dell'ingegneria.
  }

% ── RESEARCH & PROJECTS ──────────────────────────────────────
\section{PROGETTI DI RICERCA}

\wideentry
  {Suite CFD completa pensata per sviluppatori}
  {Progetto personale — OpenFOAM, Gmsh (fork), Python, C++}
  {Studi dei carichi aerodinamici aperti e di facile utilizzo, ottimizzati per problemi assialsimmetrici a numeri di $\mathcal{R}e$ molto elevati.}

\wideentry
  {Suite di analisi strutturale FEM \& verifiche secondo gli Eurocodici}
  {Progetto personale — OpenSeesPy, Python}
  {Strumenti completi di progettazione, visualizzazione e verifica strutturale.}

% ── PROFESSIONAL EXPERIENCE ──────────────────────────────────
\section{ESPERIENZA PROFESSIONALE}

\entry{febbraio 2026 \newline luglio 2026}
  {Tirocinio in ingegneria delle infrastrutture ferroviarie — TRACTEBEL FRANCE}
  {Progetto CDG Express, Saint-Ouen}
  {\begin{itemize}
    \item Verifica della conformità costruttiva del CDG Express, progetto ferroviario ad alta velocità da 2,5 miliardi di euro;
    \item Analisi strutturale e verifiche secondo gli Eurocodici di opere di ingegneria civile;
    \item Revisione di relazioni di calcolo strutturale e di rilievi as-built in cantiere;
    \item Redazione di relazioni tecniche e raccomandazioni per il committente.
  \end{itemize}}

\entry{luglio 2025 \newline gennaio 2026}
  {Tirocinio Project Manager Digital Twin — SNCF Voyageurs}
  {Industria \& soluzioni digitali, Saint-Denis}
  {\begin{itemize}
    \item Gestione di un progetto digitale da 150k€ nell'ambito della manutenzione del materiale rotabile;
    \item Coordinamento di 15 artisti 3D presso gli impianti di manutenzione di tutta la Francia;
    \item Responsabilità del progetto e reporting ai clienti interni (direttori di impianto, Direttore Industriale, ecc.);
    \item Gestione completa del processo di approvvigionamento e di gara.
  \end{itemize}}



% ── TECHNICAL SKILLS ─────────────────────────────────────────
\section{COMPETENZE TECNICHE}

\skillrow{CFD \& FEM}{OpenFOAM, StarCCM+, Abaqus, OpenSeesPy, Gmsh}
\skillrow{Programmazione}{Oltre 10 anni di esperienza con Python, C++ e ambienti UNIX}

% ── LANGUAGES ────────────────────────────────────────────────
\section{LINGUE}

\noindent\small
\textbf{Francese} — Madrelingua \quad\textbf{·}\quad
\textbf{Inglese} — Avanzato (C1) \quad\textbf{·}\quad
\textbf{Italiano} — Buono (B2) \quad\textbf{·}\quad
\textbf{Tedesco} — Buono (B2)

% ── ACTIVITIES & INVOLVEMENT ─────────────────────────────────
\section{ATTIVITÀ \& IMPEGNO ASSOCIATIVO}

\wideentry
  {Soccorritore — Croce Rossa Francese}
  {Soccorritore volontario}
  {Primo soccorso durante eventi pubblici (maratone, festival, Paris-Plage) e servizio di reperibilità con i Vigili del Fuoco di Parigi.}

\wideentry
  {Direttore — Compagnia studentesca di musical}
  {BridgeWay Comédie Musicale, Champs-sur-Marne}
  {Direzione di una compagnia di 75 artisti e di un team creativo e tecnico di 12 persone, con un budget di 4,5k€.}

\wideentry
  {Segretario Generale — Associazione studentesca}
  {Union des Élèves de l'École nationale des ponts et chaussées}
  {Gestione di un team di 15 persone per eventi professionali e sociali con oltre 1.000 partecipanti; gestione amministrativa di un'associazione senza scopo di lucro di diritto francese (Loi 1901).}

\vfill
\noindent{\scriptsize\textcolor{muted}{Autorizzo il trattamento dei miei dati personali ai sensi del D.Lgs. 196/2003 e del Regolamento UE 2016/679 (GDPR).}}

\end{document}
